\chapter{ROS containers and robots}\label{ch:containers}

ROS lives in Ubuntu containers (ROS~1 Noetic on 20.04, ROS~2 Humble on 22.04, Jazzy on 24.04) and
on the robots. The Mac has no ROS at all: tmux runs on the Mac and its panes SSH into a container
or a robot, where zsh, Neovim and the ROS helpers are installed.

\section{The container}

Share the configuration (mount the repository), never the data: plugins and Mason binaries are
built per operating system, so they live in volumes.

\begin{plain}
# compose.yaml in the ROS workspace
services:
  humble:
    build: .
    container_name: humble
    network_mode: host
    working_dir: /home/dev/ros2_ws
    environment:
      - ROS_DOMAIN_ID=${ROS_DOMAIN_ID:-0}
    volumes:
      - ./:/home/dev/ros2_ws
      - ~/dotfiles:/home/dev/dotfiles          # config, shared with the Mac
      - nvim-data:/home/dev/.local/share/nvim  # plugins + Mason (Linux binaries)
      - zsh-data:/home/dev/.local/share/zsh    # zsh plugins + history
volumes:
  nvim-data:
  zsh-data:
\end{plain}

The image needs a \texttt{dev} user and an SSH server, one port per container: 2222 humble, 2223
jazzy, 2224 noetic. With bridge networking, keep port 22 and publish it on localhost only, with
\texttt{ports: ["127.0.0.1:2222:22"]} in the compose file.

\begin{plain}
FROM ros:humble                            # or ros:jazzy, ros:noetic
RUN apt-get update && apt-get install -y openssh-server sudo && \
    useradd -m -s /bin/bash -G sudo dev && \
    echo "dev ALL=(ALL) NOPASSWD:ALL" >/etc/sudoers.d/dev && \
    mkdir -p /run/sshd && sed -i 's/^#\?Port .*/Port 2222/' /etc/ssh/sshd_config
COPY --chown=dev:dev authorized_keys /home/dev/.ssh/authorized_keys
CMD ["/usr/sbin/sshd", "-D"]
\end{plain}

Host networking on a Mac works out of the box in OrbStack; Docker Desktop needs it switched on in
its settings.

\section{Install the dotfiles inside}

Once per container; the volumes keep the result across rebuilds:

\begin{shell}
docker exec -it -u dev humble bash
~/dotfiles/container/install.sh       # 2-4 minutes, x86_64 or arm64
\end{shell}

It installs zsh, Neovim, fzf, starship, zoxide, lazygit, Node and eza, links the \texttt{zsh},
\texttt{nvim}, \texttt{starship}, \texttt{git} and \texttt{lazygit} configs, installs the plugins
and makes zsh the login shell. Options: \cmd{--no-apt} when the image already has the apt packages,
\cmd{--minimal} to skip Node (faster, but language servers that need npm will not install).

On Ubuntu 20.04 (glibc 2.31) the installer takes Neovim from \texttt{neovim/neovim-releases},
built for old glibc, and skips the tree-sitter CLI; Neovim then uses its bundled parsers.

\section{SSH into the container}

\file{install.sh} on the Mac created \file{~/.ssh/config.d/containers}:

\begin{keys}{2.4cm}
\head{Host}{Container}
\cmd{humble} & ROS~2 Humble, Ubuntu 22.04, \texttt{localhost:2222}\\
\cmd{jazzy} & ROS~2 Jazzy, Ubuntu 24.04, \texttt{localhost:2223}\\
\cmd{noetic} & ROS~1 Noetic, Ubuntu 20.04, \texttt{localhost:2224}\\
\end{keys}

All of them connect as \texttt{dev}, forward your SSH agent (so \cmd{git push} works inside) and do
not pin host keys, because containers are rebuilt often. Thanks to connection sharing, the first
connection authenticates and every further pane to the same host opens instantly.

\begin{keys}{3.6cm}
\head{Keys / command}{Action}
\cmd{ssh humble} & One shell in the container\\
\cmd{dev humble} & tmux session \texttt{humble} with windows \texttt{code}, \texttt{build}, \texttt{run}, each already in the container\\
\key{C-a S} or \cmd{s} & Pick any host\\
\key{C-a d} or \cmd{dsh} & \cmd{docker exec} shell in a running container, no SSH needed\\
\end{keys}

\section{ROS helpers}

The helpers exist only where \file{/opt/ros} exists. The installed distro is sourced
automatically (\texttt{export ROS\_AUTO=0} turns that off) and the helpers know whether they run
under ROS~1 or ROS~2.

\begin{keys}{3.4cm}
\head{Command}{Action}
\cmd{ros [distro]} & Source \file{/opt/ros/<distro>} and the nearest workspace overlay\\
\cmd{rosws} & Print the workspace root, also before the first build\\
\cmd{cb} / \cmd{cbp <pkg>} & Build the workspace / packages with their dependencies; writes \file{compile_commands.json} for clangd. colcon on ROS~2; \cmd{catkin build} or \cmd{catkin_make} on ROS~1\\
\cmd{ct [pkg]} & Run the tests and print the results\\
\cmd{rcl} & Remove \file{build/ install/ devel/ log/} (asks first)\\
\cmd{rosdomain <id>} & ROS~2: set \texttt{ROS\_DOMAIN\_ID}, shown in the prompt as \texttt{\#id}\\
\cmd{rosmaster <host>} & ROS~1: set \texttt{ROS\_MASTER\_URI} and \texttt{ROS\_IP}, shown as \texttt{@host}\\
\cmd{rtopic} / \cmd{rnode} & Pick a topic and echo it / pick a node and show its info\\
\cmd{rtl}, \cmd{rte}, \cmd{rnl} & Topic list, topic echo, node list\\
\cmd{rl}, \cmd{rr}, \cmd{rbag} & Launch, run, bag\\
\end{keys}

\section{ROS in Neovim}

The \key{Space r} menu appears only where ROS is available (a sourced distro, or colcon or catkin
on the \texttt{PATH}). Start Neovim from a shell where ROS is sourced.

\screenshot[0.62]{ros-menu}{\key{Space r}: the ROS menu in which-key, inside the \emph{humble} container.}

\begin{keys}{2.6cm}
\head{Keys}{Action}
\key{Space rb} & Build the package of the current file in the background: errors go to the quickfix list, \file{compile_commands.json} is merged and clangd restarts\\
\key{Space rB} & Build the whole workspace\\
\key{Space rL} & Last build log\\
\key{Space rc} & Merge \file{compile_commands.json}\\
\key{Space rt} & Topics: echo, hz or info\\
\key{Space rn} & Nodes: info\\
\key{Space rl} & Pick a launch file: \cmd{ros2 launch} / \cmd{roslaunch}\\
\key{Space rr} & Pick an executable: \cmd{ros2 run} / \cmd{rosrun}\\
\end{keys}

Commands: \cmd{:RosBuild [pkg]}, \cmd{:RosLog}, \cmd{:RosCompileDb}.

\screenshot[0.82]{ros-quickfix}{\key{Space rb} with a compiler error: the error is in the quickfix
list and shown on the line; \key{]q} / \key{[q} jump between errors.}

For clangd, copy \file{.clangd} and \file{.clang-format} from \file{templates/ros2/} into the
workspace root. The container installer also writes \file{~/.colcon/defaults.yaml}, so a plain
\cmd{colcon build} exports compile commands too.

\section{Robots and Jetsons}

The same installer works directly on a robot:

\begin{shell}
git clone https://github.com/agus-darmawan/dotfiles ~/dotfiles
~/dotfiles/container/install.sh --minimal
\end{shell}

On the Mac, add the robot to \file{~/.ssh/config.d/robots}, copy your key once with
\cmd{skey go2}, then:

\begin{keys}{4.2cm}
\head{Command}{Action}
\cmd{dev go2} & tmux session on the robot (windows code, build, run)\\
\cmd{sall go2 lite3} & One window, one pane per robot, typing goes to all of them (\key{C-a y} toggles)\\
\cmd{sfwd <host> <port>} & Forward a remote port (dashboard, Foxglove) to \texttt{localhost}\\
\cmd{skey <host>} & Copy your ed25519 key to a host\\
\end{keys}

Without internet on the robot, add \texttt{export NVIM\_OFFLINE=1} to its \file{~/.zshrc.local}. When
the robot runs its own tmux, press \key{F12} on the Mac (section~\ref{ch:tmux}).
